\section{Related Literature}
\label{sec:intro}

This paper contributes to three literatures: on housing costs and fertility,
on earnings risk and fertility, and on housing over the life cycle.

\paragraph{Housing costs and fertility.}
A large empirical literature finds that housing conditions shape
childbearing. Fertility is lower where rent per room is higher
\citep{simonTamuraHigherRents2009}, dwelling size predicts first births
\citep{stromHousingFirstBirths2010}, and higher house prices raise births
among owners but lower them among renters
\citep{dettlingHousePricesBirth2014,lovenheimFamilyWealthShocks2013}. A
second group of papers finds that access to mortgage credit raises births
\citep{hacamoBabiesMortgageDeregulation2021,cummingMonetaryPolicyBirth2024,dettlingKearneyModernMortgage2025,vandoornikFazioRamadoraiSkrastinsHousingFertility2025}.
Two recent papers model the joint choice of housing and children.
\citet{khorunzhinaMillerAmericanDream2022} estimate a dynamic model of
female labor supply, fertility and the first home purchase at given house
prices, and \citet{couillardBuildBabyBuild2025} estimates a model of
fertility, living arrangements and location in which rents are the cost of
housing. Relative to these papers, I study fertility together with saving,
the choice of space and tenure, and uninsurable earnings risk, in a housing
market with an equilibrium price. This separates the recurring cost of
family space from the equity needed to buy it, and allows me to compare
housing supply and mortgage credit within the same economy.

\paragraph{Earnings risk and fertility.}
A child is an irreversible commitment, and prudent parents facing more risk
want fewer children \citep{fraserIncomeRisk2001}. Quantitative life-cycle
models find that uninsurable earnings risk delays and reduces births
\citep{sommerFertilityChoiceLife2016,ejrnaesJorgensenFamilyPlanning2020,nakakuniFertilityIncomeUncertainty2026},
and \citet{deLaCroixPommeretChildbearing2021} show that risk gives the
postponement of a first birth an option value. Job displacement and job
insecurity lower fertility in quasi-experimental studies
\citep{delBonoWeberWinterEbmer2012,clarkLepinteurJobInsecurity2022}. None
of these papers includes housing. In this paper, part of the commitment a
child entails is the space it needs, which must be paid for every period at
a price set in the housing market. The fertility response to earnings risk
therefore depends on housing costs and on the down payment required to own
family space.

\paragraph{Housing over the life cycle.}
Quantitative models of housing study tenure choice under earnings risk,
down payments and transaction costs
\citep{fisherGervaisHomeOwnership2011,halketVasudevSaving2014,boarGoreaMidriganLiquidity2022},
and the joint determination of house prices, rents and ownership
\citep{kaplanMitmanViolanteHousing2020,greenwaldGurenCredit2025}. Rising
earnings risk accounts for much of the decline in homeownership among young
households
\citep{fisherGervaisHomeOwnership2011,pazPardoHomeownership2024}. In these
models, household size is exogenous or absent. I make fertility a choice. A
birth raises the household's demand for space over its life cycle and, a
generation later, the number of households competing for housing, so
housing policy shapes both family size and the future demand for housing.
